\section{Related Work}\label{sec:related}

\subsection{Monocular depth estimation}
Monocular depth models fall into two families according to what their output promises.

\paragraph{Relative / affine-invariant} MiDaS~\cite{midas} introduced training across many datasets with a loss that
ignores scale and shift, giving strong cross-domain generalisation at the price of an output that is \emph{disparity up
to an affine transform}: two constants $(s, t)$ per image must be known to obtain metres. Depth Anything~\cite{da1}
and Depth Anything~V2~\cite{da2} extend the idea with large unlabelled corpora and synthetic teachers, giving much sharper
edges and detail while remaining affine-invariant in disparity. We use DA-v2 (Large / Small) and MiDaS small as
representatives of this family and \textbf{treat $(s,t)$ as an experimental variable} (the \code{ScaleAligner}) rather
than hiding it in the model.

\paragraph{Metric} ZoeDepth~\cite{zoedepth} and the metric variants of DA-v2 (fine-tuned on Hypersim / Virtual KITTI)
aim to output metres directly. Metric3D~\cite{metric3d} identifies the camera focal length as the main ambiguity and
resolves it by normalising images to a canonical camera. These works report per-image AbsRel/$\delta_1$ on NYUv2/KITTI;
our results (\S\ref{sec:metric}) show that on real room video the per-frame scale of a metric model still swings from
$-60\%$ to $+80\%$ within one scan---invisible to per-image metrics, immediately visible once frames are fused into one
volume.

\paragraph{Conditioning on the camera, or on a teacher} Two families answer the scale ambiguity differently. One
conditions the network on camera intrinsics---Metric3D~\cite{metric3d}, UniDepth~\cite{unidepth}, and Depth
Pro~\cite{depthpro}, which estimates the focal length itself---so a metre prediction is tied to a known field of view.
The other adapts a pretrained model to the target domain with depth labels, which is what the metric DA-v2 variants do
on Hypersim. We take the second route but change where the labels come from: instead of a laser scanner or a synthetic
renderer, the teacher is the depth sensor already shipping in consumer phones. This puts a price on the teacher---the
question a product team actually faces---rather than on the architecture, and it is measured against a laser scanner
both in mesh centimetres on held-out rooms and in per-frame AbsRel on 26 rooms, not on a single-image benchmark
(\S\ref{sec:main}). The two families are complementary: an
intrinsics-conditioned model could be fine-tuned with the same LiDAR labels. We measure one of them
(Depth Pro, handed the camera's true focal length) as the Level-0 row of Exp.~6: on these captures it does
\emph{not} deliver usable metric depth (\S\ref{sec:finetune}), which is why fine-tuning is the axis we study.

\paragraph{Mobile LiDAR depth} ARKit provides $256\times192$ depth fusing LiDAR with imagery; ARKitScenes~\cite{arkitscenes}
is the first dataset to release both this depth and Faro laser-scanner depth for the same rooms, making a direct
``mobile LiDAR vs.\ laser'' comparison possible. We use both as two rows of the same table.

\subsection{Multi-view 3D reconstruction from video}
\paragraph{Classical depth fusion} Volumetric TSDF fusion~\cite{curless96} and KinectFusion~\cite{kinectfusion}
combine many depth frames through a truncated signed distance per voxel and extract the surface with marching cubes.
The property we need is that \textbf{fusion neither knows nor cares where depth comes from}---as long as it is metric
and posed---so it is a fair ``fixed part'' for every source. We use the Open3D implementation~\cite{open3d}.

\paragraph{End-to-end learned reconstruction} Atlas~\cite{atlas} regresses a TSDF directly from multi-image features,
NeuralRecon~\cite{neuralrecon} does so incrementally in real time, and SimpleRecon~\cite{simplerecon} shows that good
multi-view depth plus classical fusion is competitive. These are strong ScanNet baselines, but all learn depth
\emph{from multiple images}; they do not answer ``what room do you get from a per-image monocular model plus the
device's VIO pose,'' which is the situation of an app on a device without LiDAR. We therefore \textbf{do not compare
against this group} and make no claim of superiority.

\paragraph{Scale in monocular SLAM / reconstruction} Using monocular depth inside SLAM requires per-frame scale
recovery, e.g.\ from sparse tracked points (CNN-SLAM~\cite{cnnslam} and successors), or joint optimisation of
scale/shift with pose (MonoSDF~\cite{monosdf} uses monocular priors in a neural surface). Our \code{sparse\_points}
aligner follows this line; we evaluate it next to \code{oracle\_frame} and \code{per\_scene} so that the ``price of not
knowing scale'' is measured at three levels of available metric information.

\subsection{Room-scanning products}
Apple RoomPlan~\cite{roomplan} produces a floor plan and furniture boxes from LiDAR on iPhone/iPad~Pro; Polycam and
similar apps produce meshes or photogrammetry. All require LiDAR for room mode or fall back to slow offline
photogrammetry with many photos. No public report states how many centimetres these products deviate from a laser
scanner; this paper provides that number at least for our own pipeline.

\subsection{Position of this work}
\begin{table*}[t]
\centering\small
\begin{adjustbox}{max width=\textwidth}
\begin{tabular}{@{}lcccc@{}}
\toprule
 & Trains & Compares depth sources & Reference & Metrics \\
\midrule
Monocular depth papers & \checkmark & $\times$ (one model) & per-image depth & AbsRel, $\delta$ \\
Atlas / NeuralRecon / SimpleRecon & \checkmark & $\times$ (multi-view only) & ScanNet mesh & Chamfer, F-score \\
\textbf{This work} & head only (LiDAR teacher) & \checkmark\ laser / LiDAR / mono $\times$ aligner $\times$ teacher & Faro laser (ARKitScenes) & Chamfer, F@$\tau$, per-frame AbsRel/$\delta$, time \\
\bottomrule
\end{tabular}
\end{adjustbox}
\end{table*}
